\reading{CR-14}{Derivatives (Gregory Ch.2)}{DEFER}{1}

\begin{crkeybox}[Learning objectives for this reading]
\begin{enumerate}[label=\textbf{\alph*.},leftmargin=2em]
\item Define derivatives and explain how derivative transactions create counterparty credit risk.
\item Compare and contrast exchange-traded derivatives and over-the-counter (OTC) derivatives, and discuss the features of their markets.
\item Describe the process of clearing a derivative transaction.
\item Identify the participants and describe the use of collateralization in the derivatives market.
\item Define the International Swaps and Derivatives Association (ISDA) Master Agreement, the risk-mitigating features it provides, and the default events it covers.
\item Describe the features and use of credit derivatives and discuss potential risks they may create.
\item Describe central clearing of OTC derivatives and discuss the roles, mandate, advantages, and disadvantages of the central counterparty (CCP).
\item Explain the margin requirements for both centrally-cleared and non-centrally-cleared derivatives.
\item Define special purpose vehicles (SPVs), derivatives product companies (DPCs), monolines, and credit derivatives product companies (CDPCs) and describe the limitations of using them as risk mitigating methods.
\item Describe the approaches used and the challenges faced in modeling derivatives risk.
\end{enumerate}
\end{crkeybox}

% =====================================================================
\lo{CR-14 a}{Define derivatives and explain how derivative transactions create counterparty credit risk.}

\begin{crdefbox}[Derivative]
A derivative represents a contractual agreement between two parties to buy or
sell an underlying security, or to make a payment in the future. Derivatives can
be short-term or long-term, and are used to manage risk (hedging) or to speculate
on price movements. There are standardized versions, like futures, and custom
ones, like swaps.
\end{crdefbox}

\begin{crdefbox}[Counterparty risk]
Counterparty risk is the \emph{conditional} risk that the NPV of the investment
is positive, but the counterparty fails to perform its obligations, so no profit
is realized from the trade.
\end{crdefbox}

Counterparty risk is managed over time through \term{clearing}. This can be
performed \emph{bilaterally}, where each counterparty manages the risk of the
other, or \emph{centrally} through a central counterparty (CCP).

Products carrying credit risk include loans, forwards, swaps, options and exotic
options.

% =====================================================================
\lo{CR-14 b}{Compare and contrast exchange-traded derivatives and over-the-counter (OTC) derivatives, and discuss the features of their markets.}

\begin{center}
\footnotesize
\begin{tabular}{L{2.6cm} L{5.7cm} L{6.2cm}}
\toprule
 & \thead{Exchange-traded} & \thead{Over-the-counter (OTC)} \\
\midrule
Contract & Standardized contracts with fixed terms: price, expiry, notional
  amount & Privately negotiated between two parties, highly customizable to meet
  specific needs \\
Liquidity & Highly liquid and transparent &
  Lower liquidity; most OTC markets are less liquid, though some, like foreign
  exchange, remain highly liquid \\
Clearing and settlement & Easy to clear and settle &
  Harder to trade and unwind \\
Flexibility & Lacks flexibility, since terms cannot be customized &
  Allows precise risk management \\
Exiting a position & Straightforward &
  May require counterparty approval (\term{novation}) and can be costly or
  unfavourable \\
\bottomrule
\end{tabular}
\end{center}
\figcap{Every row is a trade-off of the same kind: standardization buys liquidity
and cheap exit, customization buys precision and pays for it in both.}

% =====================================================================
\lo{CR-14 c}{Describe the process of clearing a derivative transaction.}

\begin{crdefbox}[Clearing and settlement]
\term{Clearing} is the process between \emph{execution} and \emph{settlement} in
OTC derivative transactions. It records trade details, calculates margins, and
ensures obligations are met.

\smallskip
\term{Settlement} is the final step: the actual transfer of cash or assets to
fulfil the contract.

\smallskip
Both help reduce counterparty credit risk, that is, non-payment or non-delivery
risk.
\end{crdefbox}

\begin{center}
\begin{tikzpicture}[font=\sffamily\footnotesize]
\node[draw=FRMNavy, fill=PaleNavy, line width=0.8pt, rounded corners=1pt,
  text width=3.7cm, align=center, minimum height=1.6cm] (ex) at (0,0)
  {\textbf{EXECUTION}\\[1pt]\scriptsize buyer and seller enter a legal obligation to buy or sell};
\node[draw=FRMGold, fill=PaleGold, line width=0.8pt, rounded corners=1pt,
  text width=3.7cm, align=center, minimum height=1.6cm] (cl) at (5.2,0)
  {\textbf{CLEARING}\\[1pt]\scriptsize transaction is managed prior to settlement: margining, cash flow, payments};
\node[draw=FRMGreen, fill=PaleGreen, line width=0.8pt, rounded corners=1pt,
  text width=3.7cm, align=center, minimum height=1.6cm] (se) at (10.4,0)
  {\textbf{SETTLEMENT}\\[1pt]\scriptsize securities and/or cash exchanged; legal obligations fulfilled};
\draw[-{Stealth[length=5pt]}, FRMNavy, line width=1pt] (ex) -- (cl);
\draw[-{Stealth[length=5pt]}, FRMGold, line width=1pt] (cl) -- (se);
% routing
\node[draw=FRMRule, fill=white, line width=0.6pt, text width=2.9cm,
  align=center] (et) at (1.2,-2.6) {\scriptsize Exchange-traded};
\node[draw=FRMRule, fill=white, line width=0.6pt, text width=2.9cm,
  align=center] (otc) at (1.2,-4.1) {\scriptsize OTC-traded};
\node[draw=FRMGreen, fill=PaleGreen, line width=0.7pt, text width=3.4cm,
  align=center] (cc1) at (6.6,-2.6) {\scriptsize centrally cleared};
\node[draw=FRMGreen, fill=PaleGreen, line width=0.7pt, text width=3.4cm,
  align=center] (cc2) at (6.6,-3.75) {\scriptsize centrally cleared};
\node[draw=FRMRed, fill=PaleRed, line width=0.7pt, text width=3.4cm,
  align=center] (bc) at (6.6,-4.75) {\scriptsize bilaterally cleared};
\draw[-{Stealth[length=4.5pt]}, FRMGrey, line width=0.8pt] (et) -- (cc1);
\draw[-{Stealth[length=4.5pt]}, FRMGrey, line width=0.8pt] (otc) -- (cc2);
\draw[-{Stealth[length=4.5pt]}, FRMGrey, line width=0.8pt] (otc) -- (bc);
\end{tikzpicture}
\end{center}
\figcap{Clearing sits between execution and settlement, not at either end. The
lower panel is the fact that gets tested: an exchange-traded derivative has only
one route, and it is central clearing, while an OTC derivative can go either way.}

\begin{itemize}
\item \term{Bilateral clearing.} Counterparties manage risk independently. This is
the traditional OTC market, and it remains prevalent.
\item \term{Central clearing.} A third party, often a CCP, manages counterparty
risk and collateralization. It requires standardization to implement and is not
universally applicable to all OTC derivatives. \term{Multilateral compression}
mimics central clearing functionalities.
\end{itemize}

% =====================================================================
\lo{CR-14 d}{Identify the participants and describe the use of collateralization in the derivatives market.}

\subsection*{Participants, by portfolio size}

\begin{enumerate}
\item \term{Large players (dealers): global banks.} Trade most derivatives and
provide liquidity. Engage in extensive derivatives trading. Members of exchanges
and central clearinghouses.
\item \term{Medium-sized players: smaller banks and other financial institutions}
such as hedge funds and pension funds. Involved in OTC derivatives activities;
members of local exchanges and clearinghouses.
\item \term{End users: large corporations, sovereigns, smaller financial
institutions.} Use derivatives for hedging or investment, with limited
derivatives use.
\item \term{Third-party service providers.} Offer settlement, margining,
collateral management, software, trade compression, and clearing services.
\end{enumerate}

\subsection*{Collateralization, by derivative type}

\begin{center}
\begin{tikzpicture}[font=\sffamily\footnotesize]
\foreach \i/\lab/\det/\col in {
  0/{Exchange traded}/{simple, liquid, short-dated; centrally cleared; daily margin posting}/FRMGreen,
  1/{OTC centrally cleared}/{illiquid or non-standardized, but centrally cleared; daily variation margin in cash}/FRMGreen,
  2/{OTC collateralized}/{bilateral, no central clearing; parties post cash or securities}/FRMGold,
  3/{OTC uncollateralized}/{bilateral; no collateral posted, or less and lower-quality collateral}/FRMRed}
{
  \node[draw=\col, fill=\col!10, line width=0.8pt, rounded corners=1pt,
    text width=3.2cm, align=center, minimum height=2.0cm] at ({\i*3.75},0)
    {\scriptsize\textbf{\lab}\\[2pt]\scriptsize\det};
}
\draw[-{Stealth[length=6pt]}, FRMRed, line width=1pt] (-1.85,-1.45) -- (13.1,-1.45)
  node[midway, below, font=\scriptsize\itshape, text=FRMRed, fill=white,
  inner sep=1.5pt] {increasing counterparty risk};
\end{tikzpicture}
\end{center}
\figcap{The four tiers, ordered by the counterparty risk left standing. What moves
along the row is not the product but the \emph{protection}: central clearing
first, then daily margin, then collateral of any kind, then nothing.}

% =====================================================================
\lo{CR-14 e}{Define the International Swaps and Derivatives Association (ISDA) Master Agreement, the risk-mitigating features it provides, and the default events it covers.}

\begin{crdefbox}[ISDA Master Agreement]
Established by the International Swaps and Derivatives Association in 1985 and
widely adopted since 1992, the ISDA Master Agreement is the industry standard
legal framework governing OTC derivative transactions. Multiple transactions are
covered under a general Master Agreement to form a \term{single legal contract} of
indefinite term. Its components are a common core and an adjustable schedule for
flexibility. Its objectives are to eliminate legal uncertainties and mitigate
counterparty risk.
\end{crdefbox}

\begin{center}
\footnotesize
\begin{tabular}{L{7.0cm} L{7.5cm}}
\toprule
\thead{Risk-mitigating features} & \thead{Default events covered} \\
\midrule
Collateral terms & Failure to pay or deliver \\
Default events and termination & Breach of agreement \\
Netting of transactions & Credit support default \\
A defined close-out process mechanism & Misrepresentation \\
 & Default under specified transaction \\
 & Cross-default \\
 & Bankruptcy \\
 & Merger without assumption \\
\bottomrule
\end{tabular}
\end{center}
\figcap{Four risk-mitigating features and eight default events. The objective
names both halves explicitly, and the eight events are the half most often
skipped.}

% =====================================================================
\lo{CR-14 f}{Describe the features and use of credit derivatives and discuss potential risks they may create.}

Credit derivatives are complex financial instruments that can be used to hedge
credit risk, but they introduce new risks. The CDS is the most common type: the
buyer pays a premium for protection against a specific borrower's default.

\begin{crtrapbox}[The risk they create]
The credit derivatives market can be susceptible to \term{speculation}, which can
increase risk in the financial system. The instrument that exists to transfer
credit risk can therefore also be used to take it on by parties with no
underlying exposure.
\end{crtrapbox}

% =====================================================================
\lo{CR-14 g}{Describe central clearing of OTC derivatives and discuss the roles, mandate, advantages, and disadvantages of the central counterparty (CCP).}

\subsection*{Role of the CCP}

\begin{itemize}
\item Interposes itself between counterparties, acting as \emph{buyer to sellers
and seller to buyers}.
\item Reallocates default losses through \term{netting, margining, and loss
mutualisation}.
\item Sets standards and rules for clearing members.
\item Closes out positions of defaulting clearing members.
\item Maintains financial resources --- cash variation margin, initial margin,
default fund --- to cover losses.
\item Has a documented plan for the very extreme situation in which all its
financial resources are depleted.
\end{itemize}

\subsection*{The clearing mandate after the crisis}

Prior to the global financial crisis there was no clear requirement for central
clearing of OTC derivatives. Calls for central clearing emerged globally after
Bear Stearns. \term{G20 leaders agreed in Pittsburgh in 2009} on requirements
including:

\begin{itemize}
\item all standardised OTC derivatives to be traded on exchanges or electronic
platforms;
\item mandatory central clearing of standardised OTC derivatives;
\item reporting of OTC derivatives to trade repositories; and
\item higher capital requirements for non-centrally cleared OTC derivatives.
\end{itemize}

\begin{center}
\footnotesize
\begin{tabular}{L{7.0cm} L{7.5cm}}
\toprule
\thead{Advantages of central clearing} & \thead{Disadvantages of central clearing} \\
\midrule
\term{Netting}: reduces counterparty risk by allowing netting of all trades &
  \term{Loss mutualisation} may disperse losses across all clearing members \\
\term{Margin management}: manages risks associated with portfolio value
  fluctuations &
  \term{Large price moves} and increased volatilities during close-outs can
  result in market impact \\
\term{Financial shock absorber}: absorbs the domino effect in case of member
  default, ensuring swift termination of financial relations &
  Requires sufficient financial resources, potentially burdening members \\
 & Not all OTC derivatives can be centrally cleared, due to standardization and
  liquidity requirements \\
\bottomrule
\end{tabular}
\end{center}

\begin{crtrapbox}[Netting appears on both sides of the ledger]
Netting is listed as an advantage, and loss mutualisation --- the mechanism that
makes CCP netting work --- is listed as a disadvantage. They are the same
property viewed from two positions: netting is good for the system and loss
mutualisation is bad for the surviving members who fund it.
\end{crtrapbox}

Two further shortfalls: CCPs must be \emph{highly creditworthy} to prevent
systemic shocks, and their cross-border activities are complicated by varying
legal and regulatory environments. Reducing derivatives risk may also disadvantage
other market participants, such as bondholders.

% =====================================================================
\lo{CR-14 h}{Explain the margin requirements for both centrally-cleared and non-centrally-cleared derivatives.}

\begin{crkeybox}[The two regimes]
\term{Centrally cleared.} Clearing members must post \emph{initial margin} and
\emph{variation margin} and contribute to a \emph{default fund} to cover the
losses of a defaulting clearing member. Initial margin is required on all trades;
variation margin covers market value changes in the derivatives.

\smallskip
\term{Non-centrally cleared (bilateral).} Historically had no or only limited
margin requirements. In recent years, to mitigate systemic risk, regulators
globally have been moving to introduce margin requirements for bilateral trades
similar to those under central clearing.
\end{crkeybox}

% =====================================================================
\lo{CR-14 i}{Define special purpose vehicles (SPVs), derivatives product companies (DPCs), monolines, and credit derivatives product companies (CDPCs) and describe the limitations of using them as risk mitigating methods.}

\subsection*{Special purpose vehicle (SPV)}

An SPV, also known as a special purpose entity, is an off-balance sheet,
\term{bankruptcy-remote} entity that is legally separate from its sponsor. It is
created primarily to reduce counterparty credit risk by isolating assets and
liabilities from the parent institution, and is frequently used in structured
notes.

The structure works by transferring assets from the sponsor to the SPV so that
investors in the SPV have priority claims, a priority often called a \term{flip
provision}.

\begin{crtrapbox}[SPVs transform risk rather than eliminate it]
SPVs do not eliminate risk; they transform \emph{counterparty} risk into
\emph{legal} risk. In the case of sponsor default, a bankruptcy court may
consolidate the SPV's assets with the sponsor, treating them as if they were never
transferred, which defeats the purpose of the isolation.

\smallskip
Jurisdiction matters: U.S.\ courts are more likely to allow consolidation, while
U.K.\ courts avoid consolidation unless there is fraud. In the Lehman case a flip
provision was declared unenforceable in the U.S., contradicting a U.K.\ ruling,
which highlights the legal uncertainty.
\end{crtrapbox}

\subsection*{Derivatives product company (DPC)}

A DPC is a separately capitalized entity created by a financial institution to
achieve a high (AAA) credit rating and reduce counterparty risk in OTC
derivatives. An originating bank sets up a bankruptcy-remote SPV and injects
capital to gain a preferential and strong credit rating, typically triple-A. It
works by distancing itself from the parent so that it can meet obligations even if
the parent fails.

Its credit quality depends on three factors:

\begin{enumerate}[label=\alph*)]
\item minimizing market risk through offsetting trades (\term{mirrored trades});
\item structural support from the parent (transfer or orderly unwind); and
\item strong internal risk management (mark-to-market and collateral).
\end{enumerate}

During the global financial crisis these assumptions failed. Lehman's DPCs went
bankrupt and Bear Stearns' DPCs were stressed, showing that DPCs still face legal,
market and operational risks.

\subsection*{Monolines}

Financial guarantee insurance companies with a \emph{single} business line,
primarily providing credit enhancement (\term{credit wraps}) to improve bond
ratings, especially in structured credit markets. They typically maintain very
high credit ratings (AAA) and are active in CDS and structured products, allowing
issuers to borrow at lower costs.

\begin{crtrapbox}[The monoline failure mechanism]
Monolines \term{do not post collateral} in normal market conditions, relying
instead on their strong credit rating. During the crisis, losses on insured
products rose, ratings were downgraded, and the downgrades \emph{triggered
requirements to post collateral} --- at exactly the moment insured losses were
rising. Several monoline insurers failed between 2007 and 2009.
\end{crtrapbox}

\subsection*{Credit derivative product companies (CDPCs)}

CDPCs are similar to DPCs in structure but operate like monolines. They provide
credit protection, often through instruments similar to credit default swaps, and
generate income from the premiums received.

\begin{itemize}
\item Unlike DPCs, CDPCs \emph{do not} typically hedge their exposures with
offsetting positions, so they are \textbf{not market neutral}.
\item They are usually highly leveraged and retain significant exposure to credit
risk.
\item Like monolines, they do not post collateral in normal conditions and rely on
maintaining high credit ratings.
\end{itemize}

During and after the crisis many CDPCs experienced large losses and failed.

\begin{center}
\footnotesize
\begin{tabular}{L{2.2cm} L{4.0cm} L{4.0cm} L{3.6cm}}
\toprule
 & \thead{Hedged with offsetting trades?} & \thead{Posts collateral in normal conditions?} &
\thead{Principal weakness} \\
\midrule
\term{SPV} & --- & --- & Legal risk: court consolidation \\
\term{DPC} & \textbf{Yes} (mirrored trades) & Yes (mark-to-market and collateral) &
  Depends on parent support that failed in the crisis \\
\term{Monoline} & No & \textbf{No} & Downgrade triggers collateral calls just as
  losses arrive \\
\term{CDPC} & \textbf{No} & \textbf{No} & Leverage plus no hedge plus reliance on
  ratings \\
\bottomrule
\end{tabular}
\end{center}
\figcap{The DPC is the odd one out: it is the only one of the four that both
hedges and posts collateral. Monolines and CDPCs are distinguished from each other
by the hedging column, not the collateral column.}

% =====================================================================
\lo{CR-14 j}{Describe the approaches used and the challenges faced in modeling derivatives risk.}

\begin{enumerate}[label=\alph*)]
\item \term{Value at risk (VaR)} measures the worst loss over a specified short
period at a given confidence level. VaR is simple and intuitive and does not
depend on the shape of the loss distribution, but it does not reveal anything
about the possible loss \emph{beyond} the confidence level.
\item \term{Potential future exposure (PFE)} is the credit risk equivalent of VaR,
relating to \emph{gains} and using a \emph{long} time horizon.
\item \term{Expected shortfall (ES)} measures the average loss knowing that the
loss is at least equal to VaR.
\end{enumerate}

\begin{crtrapbox}[The challenge]
Financial institutions often combine these model types with other complex
quantitative models using correlations, volatility and dependence factors.
However, \emph{historical correlations are often poor predictors of future
outcomes}, and they can change significantly during periods of stress --- which is
precisely when the models are being relied upon.
\end{crtrapbox}
